% =====================================================================
%  CHAPITRES 5, 6 ET 7 — Les sons, le rythme, la clinique francophone
% =====================================================================
\chapter{Les 44 sons de l'anglais}

\section{Pourquoi commencer par les sons ?}
Trois raisons, toutes décisives :

\begin{enumerate}
  \item \textbf{L'anglais n'est pas une langue phonétique.} En français, une lettre = un son,
  ou presque. En anglais, \en{ough} se prononce de six façons différentes :
  \en{through} \pron{throu}, \en{though} \pron{dzo}, \en{thought} \pron{thoot},
  \en{tough} \pron{teuf}, \en{cough} \pron{kof}, \en{bough} \pron{baou}. Si vous lisez
  « à la française », vous prononcez faux --- et vous ne reconnaissez pas le mot quand
  quelqu'un d'autre le dit.
  \item \textbf{L'anglais possède 9 sons qui n'existent pas en français.} Tant que votre
  oreille ne les distingue pas, vous entendez \en{ship} et \en{sheep} comme le même mot,
  et vous produisez \en{think} comme \en{sink}. Ce n'est pas un problème de gorge, c'est un
  problème d'\textbf{oreille} : on ne peut pas produire un son qu'on n'entend pas.
  \item \textbf{C'est le premier obstacle ressenti à l'oral.} Beaucoup de francophones
  connaissent les mots mais n'osent pas parler, parce qu'ils « sonnent faux ». Deux semaines
  de travail phonétique suppriment 80\,\% de cette gêne.
\end{enumerate}

\begin{astuce}[Comment lire les tableaux qui suivent]
Chaque son est donné sous trois formes :
\textbf{(1)} le symbole phonétique international (API) entre barres obliques ;
\textbf{(2)} un \textbf{mot anglais modèle} --- apprenez-le, il vaut mille explications ;
\textbf{(3)} une \textbf{prononciation approchée en lettres françaises}, en violet, entre
crochets. La respelling française est une béquille : elle vous dépanne les premiers mois,
mais elle est forcément imprécise. Dès que vous pouvez, remplacez-la par l'écoute réelle du
mot (Forvo, dictionnaire Cambridge avec haut-parleur, YouGlish).
\end{astuce}

\section{Les 12 voyelles simples}
C'est ici que les francophones perdent le plus. Le français a 12 à 16 voyelles selon les
analyses, l'anglais en a 12 --- mais \textbf{différentes}, et surtout l'anglais oppose des
voyelles \textbf{longues} et \textbf{brèves} qui changent le sens des mots.

\begin{tableau}[Les 12 voyelles simples (monohtongues)]
\footnotesize
\begin{longtable}{@{}p{1.7cm}p{3.6cm}p{2.6cm}p{7.2cm}@{}}
\entete{API} & \entete{Mot modèle (et autres exemples)} & \entete{À la française} & \entete{Comment la produire / piège francophone} \tabularnewline
\endfirsthead
\entete{API} & \entete{Mot modèle} & \entete{À la française} & \entete{Comment la produire / piège francophone} \tabularnewline
\endhead
\hline
/\ipaee\ipalong/ & \en{sheep}, \en{see}, \en{tea}, \en{free}, \en{easy} & \pron{chiiip} &
Long « i » tendu, lèvres étirées, son tenu. \textbf{Tenez-le} : \en{sheeeep}. Le français ne
tient pas ses voyelles --- c'est la première chose à changer. \tabularnewline
\hline
/\ipaiot/ & \en{ship}, \en{sit}, \en{big}, \en{English}, \en{minute} & \pron{chip} &
Bref, relâché, entre le « i » et le « é ». Bouche presque fermée mais muscles mous.
\textbf{La paire \en{sheep}/\en{ship} est le test classique.} \tabularnewline
\hline
/\ipaE/ & \en{bed}, \en{ten}, \en{said}, \en{many}, \en{friend} & \pron{bèd} &
« è » bref. Attention : \en{many}, \en{any}, \en{said}, \en{friend} ne se prononcent pas
comme ils s'écrivent. \tabularnewline
\hline
/\ipae/ & \en{cat}, \en{bad}, \en{man}, \en{apple}, \en{happy} & \pron{katt} &
\textbf{N'existe pas en français.} Ouvrez grand la bouche, mâchoire descendue, comme chez le
dentiste, et dites un « a » en tirant vers le « è ». Son très ouvert et nasal. \tabularnewline
\hline
/\ipaa\ipalong/ & \en{father}, \en{car}, \en{start}, \en{class} & \pron{faa-dheu} &
« a » long, arrière de la bouche, lèvres neutres. \en{AmE} \en{class} = \pron{klass},
\en{BrE} = \pron{klaas}. \tabularnewline
\hline
/\ipao/ (bref) & \en{hot}, \en{dog}, \en{stop}, \en{box} & \pron{hot} &
« o » bref et ouvert, lèvres arrondies. Américain : souvent plus proche de « aa »
(\en{hot} \pron{haat}). \tabularnewline
\hline
/\ipao\ipalong/ & \en{four}, \en{thought}, \en{law}, \en{door}, \en{more} & \pron{foo} &
« ô » long, lèvres bien arrondies et projetées. \tabularnewline
\hline
/\ipaups/ & \en{book}, \en{put}, \en{good}, \en{could}, \en{woman} & \pron{bouk} &
« ou » très bref et relâché. Ne le confondez jamais avec le « ou » long ci-dessous :
\en{full} \pron{foul} $\neq$ \en{fool} \pron{fouul}. \tabularnewline
\hline
/\ipaoo\ipalong/ & \en{too}, \en{food}, \en{blue}, \en{group}, \en{you} & \pron{touu} &
« ou » long, lèvres très arrondies, son tenu. \tabularnewline
\hline
/\ipalv/ & \en{cup}, \en{love}, \en{money}, \en{country}, \en{sun} & \pron{keup} &
\textbf{N'existe pas en français.} Bref, central, entre « a » et « eu », mâchoire relâchée.
C'est le \en{u} de \en{but}. Les francophones disent « keup » à la française : trop long et
trop rond. Soyez bref et mou. \tabularnewline
\hline
/\ipaer\ipalong/ & \en{bird}, \en{her}, \en{work}, \en{learn}, \en{early} & \pron{beuud} &
« eu » long + \textbf{r} anglais prononcé (contrairement au français où le « euh » final
disparaît). \en{bird}, \en{word}, \en{world}, \en{girl}, \en{heard} : cinq orthographes,
un seul son. \tabularnewline
\hline
/\ipeschwa/ & \en{about}, \en{teacher}, \en{banana}, \en{problem} & \pron{e} très faible &
\textbf{Le son le plus fréquent de l'anglais (environ 30\,\% des voyelles).} Voyelle neutre,
à peine prononcée, dans toutes les syllabes \emph{non accentuées}. \en{about} = \pron{e-baout},
\en{banana} = \pron{be-na-neu}, \en{teacher} = \pron{tii-tcheu}. \tabularnewline
\end{longtable}
\end{tableau}

\begin{piege}[Le français n'a pas de voyelles longues --- l'anglais si]
\en{sheep} et \en{ship}, \en{fool} et \en{full}, \en{cart} et \en{cut}, \en{beach} et
\en{bitch} : quatre paires où la longueur change le sens. Un francophone qui abrège toutes
ses voyelles produit des contresens --- parfois embarrassants. Entraînement : tenez la
voyelle longue une seconde entière, puis dites la brève en un éclair.
\end{piege}

\section{Les 8 diphtongues}
Une diphtongue est un glissement d'une voyelle vers une autre, \emph{dans la même syllabe}.
Le français en a peu (\en{oui}, \en{bien}) ; l'anglais en est plein. Le secret : le premier
élément est fort et long, le second est faible et court.

\begin{tableau}[Les 8 diphtongues]
\footnotesize
\begin{longtable}{@{}p{1.9cm}p{4.3cm}p{2.6cm}p{6.3cm}@{}}
\entete{API} & \entete{Mots modèles} & \entete{À la française} & \entete{Remarque} \tabularnewline
\endfirsthead
\entete{API} & \entete{Mots modèles} & \entete{À la française} & \entete{Remarque} \tabularnewline
\endhead
\hline
/\ipaei/ & \en{say}, \en{day}, \en{eight}, \en{name}, \en{they} & \pron{dé-i} &
Le « ai » français est \emph{une} voyelle ; l'anglais glisse. \en{day} = \pron{dé-i},
jamais \pron{dè}. \tabularnewline
\hline
/\ipaai/ & \en{my}, \en{time}, \en{five}, \en{eye}, \en{idea} & \pron{ma-i} &
Glissement net « a \ra i ». \en{time} = \pron{ta-im}, pas \pron{taïm} mou. \tabularnewline
\hline
/\ipaOi/ & \en{boy}, \en{oil}, \en{voice}, \en{enjoy} & \pron{bo-i} &
Identique au français « oi », mais plus appuyé. \tabularnewline
\hline
/\ipaau/ & \en{now}, \en{house}, \en{out}, \en{about} & \pron{na-ou} &
Commence bouche grande ouverte. \en{now} = \pron{na-ou}, pas \pron{no}. \tabularnewline
\hline
/\ipaou/ \en{(BrE)} / \ipaoou \en{(AmE)} & \en{go}, \en{home}, \en{no}, \en{show} & \pron{gou-ou} &
\textbf{Erreur n°1 des francophones :} dire \en{no} = \pron{no}. C'est \pron{nou-ou}.
Le « o » anglais est toujours une diphtongue. \tabularnewline
\hline
/\ipaiot\ipeschwa/ & \en{here}, \en{near}, \en{idea}, \en{ear} & \pron{hi-eu} &
\en{here} = \pron{hi-eu}, pas \pron{hèr}. \tabularnewline
\hline
/\ipaE\ipeschwa/ & \en{hair}, \en{care}, \en{where}, \en{there} & \pron{hè-eu} &
\en{hair} = \pron{hè-eu}. Attention : \en{here} et \en{hair} sont différents. \tabularnewline
\hline
/\ipaups\ipeschwa/ & \en{tour}, \en{pure}, \en{sure}, \en{cure} & \pron{tou-eu} &
De plus en plus prononcé \pron{too} (\en{sure} = \pron{choo}). \tabularnewline
\end{longtable}
\end{tableau}

\section{Les 24 consonnes}
Bonne nouvelle : 19 consonnes anglaises existent en français et ne posent aucun problème
(\en{p, b, t, d, k, g, f, v, s, z, m, n, l} --- avec quelques nuances). Cinq demandent un
vrai travail, signalées par \etoile.

\begin{tableau}[Les 24 consonnes]
\footnotesize
\begin{longtable}{@{}p{1.8cm}p{4.0cm}p{2.4cm}p{6.9cm}@{}}
\entete{API} & \entete{Mots modèles} & \entete{À la française} & \entete{Remarque pour francophone} \tabularnewline
\endfirsthead
\entete{API} & \entete{Mots modèles} & \entete{À la française} & \entete{Remarque pour francophone} \tabularnewline
\endhead
\hline
/p/ & \en{pen}, \en{top}, \en{apple} & \pron{pèn} &
Comme en français, mais \textbf{aspiré} en début de mot : un petit souffle
(\en{pen} = p\,+\,souffle). Devant \en{s}, pas d'aspiration : \en{spin} = \pron{spine}. \tabularnewline
\hline
/b/ & \en{bad}, \en{job}, \en{about} & \pron{bad} & Identique au français. \tabularnewline
\hline
/t/ & \en{tea}, \en{stop}, \en{water} & \pron{tii} &
Aspiré en début de mot. Américain : entre deux voyelles, il devient un « d » bref
(\en{water} = \pron{ouo-deu}, \en{city} = \pron{si-di}). \tabularnewline
\hline
/d/ & \en{day}, \en{did}, \en{red} & \pron{dé-i} & Identique. \tabularnewline
\hline
/k/ & \en{cat}, \en{school}, \en{back} & \pron{katt} &
Aspiré en début de mot. \en{school} = \pron{skoul} (pas de souffle après \en{s}). \tabularnewline
\hline
/g/ & \en{go}, \en{big}, \en{again} & \pron{gou} &
Toujours dur, jamais « j » : \en{get} = \pron{guèt}, \en{girl} = \pron{gueul}. \tabularnewline
\hline
/f/ & \en{fat}, \en{life}, \en{phone} & \pron{fatt} &
\en{ph} = \en{f} : \en{phone}, \en{photo}, \en{phrase}. \tabularnewline
\hline
/v/ & \en{van}, \en{love}, \en{very} & \pron{van} &
Ne confondez jamais \en{v} et \en{w} (voir plus bas). \tabularnewline
\hline
\etoile /\ipatheta/ & \en{think}, \en{three}, \en{thing}, \en{math} & \pron{thi-nk} &
\textbf{N'existe pas en français.} Langue \emph{entre les dents}, souffle. Jamais « s »,
jamais « f ». \en{think} $\neq$ \en{sink}. \tabularnewline
\hline
\etoile /\ipadh/ & \en{this}, \en{that}, \en{mother}, \en{the} & \pron{dhis} &
Même position de la langue, mais \textbf{avec la voix} (cordes vocales qui vibrent).
\en{this} $\neq$ \en{dis}. Ce son apparaît dans les mots les plus fréquents de l'anglais
(\en{the, this, that, they, there, them, then, than, with, mother, father, other}). \tabularnewline
\hline
/s/ & \en{see}, \en{city}, \en{miss} & \pron{sii} &
\en{c} devant \en{e/i} = \en{s} : \en{city}, \en{nice}. \tabularnewline
\hline
/z/ & \en{zoo}, \en{is}, \en{music}, \en{was} & \pron{zou} &
\en{s} entre voyelles = \en{z} : \en{rose}, \en{busy}, \en{music}. \tabularnewline
\hline
/\ipash/ & \en{she}, \en{nation}, \en{sure}, \en{special} & \pron{chii} &
Comme le « ch » français. \en{ti}, \en{ci}, \en{ss} peuvent donner ce son :
\en{nation} \pron{né-i-cheun}, \en{special} \pron{spé-chl}, \en{pressure} \pron{pré-cheu}. \tabularnewline
\hline
/\ipazh/ & \en{measure}, \en{vision}, \en{usual}, \en{pleasure} & \pron{mé-jeu} &
Comme le « j » français. Rare (environ 150 mots). \tabularnewline
\hline
\etoile /h/ & \en{hat}, \en{hello}, \en{house}, \en{behind} & \pron{hètt} &
\textbf{N'existe pas en français} (notre « h » est muet). C'est un souffle, comme sur des
lunettes pour les nettoyer. \textbf{Ne l'avalez jamais} : \en{he, him, her, house, happy,
hand, head, hold, hope}. Et ne l'ajoutez jamais là où il n'y en a pas : \en{hour},
\en{honest}, \en{honor}, \en{heir} sont muets. \tabularnewline
\hline
/t\ipash/ & \en{chair}, \en{watch}, \en{teach}, \en{nature} & \pron{tchè-eu} &
Un seul son, pas « t » + « ch » séparés. \tabularnewline
\hline
/d\ipazh/ & \en{joy}, \en{age}, \en{judge}, \en{gym} & \pron{djo-i} &
Un seul son. \en{j} = \en{d\ipazh} presque toujours (\en{jam}, \en{just}, \en{job}). \tabularnewline
\hline
/m/ & \en{man}, \en{time}, \en{come} & \pron{mann} & Identique. \tabularnewline
\hline
/n/ & \en{no}, \en{ten}, \en{when} & \pron{nou} & Identique. \tabularnewline
\hline
\etoile /\ipang/ & \en{sing}, \en{long}, \en{thing}, \en{English} & \pron{si-ng} &
\textbf{N'existe pas en français en fin de mot.} Arrière de la langue contre le palais,
\emph{sans} prononcer un « g » final. \en{sing} = \pron{si-ng} (pas \pron{si-ngue}).
Mais \en{finger} = \pron{fi-ngueu} (le \en{g} se fait entendre), \en{singer} = \pron{si-ngheu}. \tabularnewline
\hline
/l/ & \en{leg}, \en{light}, \en{all}, \en{feel} & \pron{lèg} &
Deux \en{l} anglais : clair devant voyelle (\en{light}), « sombre » --- langue en arrière ---
en fin de mot (\en{all}, \en{feel}, \en{milk}). Ne l'oubliez jamais en fin de mot. \tabularnewline
\hline
\etoile /r/ & \en{red}, \en{right}, \en{very}, \en{around} & \pron{rèd} &
\textbf{Très différent du \en{r} français} (qui vient de la gorge). Le \en{r} anglais :
langue recroquevillée \emph{au milieu} de la bouche, sans toucher le palais, lèvres un peu
arrondies. Aucun grattement. Et il se prononce \textbf{partout}, y compris avant consonne
(\en{car} \pron{kaa-r}, \en{four} \pron{foo-r}). \tabularnewline
\hline
\etoile /w/ & \en{we}, \en{water}, \en{work}, \en{want}, \en{now} & \pron{ouii} &
\textbf{N'existe pas comme phonème en français.} Lèvres très arrondies puis ouverture
rapide : c'est un « ou » qui glisse vers la voyelle suivante. \textbf{Jamais « v »} :
\en{we} $\neq$ \en{ve}, \en{wine} $\neq$ \en{vine}, \en{west} $\neq$ \en{vest}. \tabularnewline
\hline
/j/ & \en{yes}, \en{you}, \en{year}, \en{use} (verbe) & \pron{ièss} &
Le « y » anglais = « i » français de \en{bien}. Attention : \en{you} = \pron{iou},
\en{year} = \pron{ii-eu}. \tabularnewline
\end{longtable}
\end{tableau}

\begin{memo}[Les 6 sons à travailler en priorité absolue]
\begin{enumerate}
  \item \ipatheta\ et \ipadh\ (\en{think}, \en{this}) --- langue entre les dents.
  \item /h/ aspiré (\en{he}, \en{house}) --- ne jamais l'avaler.
  \item /w/ (\en{we}, \en{water}) --- ne jamais le dire « v ».
  \item /r/ anglais (\en{red}, \en{car}) --- jamais le \en{r} de gorge français.
  \item /\ipaiot/ contre /\ipaee\ipalong/ (\en{ship}/\en{sheep}) --- bref contre long.
  \item /\ipeschwa/ (le schwa) --- la clé du rythme anglais (chapitre 6).
\end{enumerate}
\end{memo}

\section{L'orthographe anglaise : ce qui se prononce et ce qui ne se prononce pas}

\subsection{Les lettres muettes (environ 60\,\% des mots anglais en contiennent)}
\begin{tableau}[Lettres muettes --- à mémoriser par groupes]
\footnotesize
\begin{longtable}{@{}p{2.2cm}p{5.2cm}p{7.2cm}@{}}
\entete{Lettre muette} & \entete{Mots} & \entete{Prononciation} \tabularnewline
\endfirsthead
\entete{Lettre muette} & \entete{Mots} & \entete{Prononciation} \tabularnewline
\endhead
\hline
\en{k} devant \en{n} & \en{know, knee, knife, knock, knight, knot} & \pron{nou}, \pron{nii}, \pron{naif}, \pron{nok}, \pron{naït} \tabularnewline
\en{w} devant \en{r} & \en{write, wrong, wrist, wrap, wreck} & \pron{raït}, \pron{rong} \tabularnewline
\en{b} final après \en{m} & \en{climb, bomb, thumb, comb, dumb, lamb} & \pron{klaïm}, \pron{bom}, \pron{theum} \tabularnewline
\en{t} dans certains mots & \en{listen, often, castle, whistle, Christmas} & \pron{li-s'n}, \pron{o-f'n}, \pron{kaa-s'l} \tabularnewline
\en{l} dans certains mots & \en{walk, talk, half, calm, would, should, could} & \pron{ouook}, \pron{toook}, \pron{haaf}, \pron{kaam}, \pron{woud} \tabularnewline
\en{g} devant \en{n} & \en{sign, design, foreign, gnaw} & \pron{saïn}, \pron{di-zaïn}, \pron{fo-rin} \tabularnewline
\en{h} initial & \en{hour, honest, honor, heir, herb (AmE)} & \pron{a-ou-eu}, \pron{o-nèst} \tabularnewline
\en{u} après \en{g} & \en{guess, guitar, guide, guest, build} & \pron{guèss}, \pron{gi-taa}, \pron{bild} \tabularnewline
\en{i} dans \en{-tion} & \en{nation, station, question, information} & \pron{né-i-cheun}, \en{question} = \pron{kwès-tchoun} \tabularnewline
\en{e} final & \en{make, time, come, house, name} & Le \en{e} est muet mais \textbf{allonge la voyelle} : \en{make} = \pron{mé-ik} et non \pron{mak}. \tabularnewline
\end{longtable}
\end{tableau}

\subsection{Un son, plusieurs orthographes}
\begin{itemize}
  \item \textbf{/\ipaee\ipalong/} : \en{see, sea, key, machine, people, believe, receive, field, these}
  \item \textbf{/\ipash/} : \en{she, nation, special, sure, ocean, conscience, sugar}
  \item \textbf{/t\ipash/} : \en{chair, watch, teach, nature, question, cello}
  \item \textbf{/\ipalv/} : \en{cup, love, money, does, blood, other, brother, enough}
  \item \textbf{/\ipaer\ipalong/} : \en{bird, work, learn, heard, girl, journey, early}
  \item \textbf{/\ipaei/} : \en{say, day, eight, name, they, rain, great, weigh, baby}
\end{itemize}

\subsection{Même orthographe, sons différents (les \en{heteronyms})}
\begin{tableau}[Mots qui changent de prononciation selon le sens]
\footnotesize
\begin{longtable}{@{}p{3.0cm}p{5.4cm}p{5.4cm}@{}}
\entete{Mot} & \entete{Nom (accent sur la 1\textsuperscript{re} syllabe)} & \entete{Verbe (accent sur la 2\textsuperscript{e} syllabe)} \tabularnewline
\endfirsthead
\entete{Mot} & \entete{Nom} & \entete{Verbe} \tabularnewline
\endhead
\hline
\en{record} & \pron{rè-kood} un disque, un record & \pron{ri-kood} enregistrer \tabularnewline
\en{present} & \pron{prè-z'nt} un cadeau & \pron{pri-zènt} présenter \tabularnewline
\en{object} & \pron{o-bdject} un objet & \pron{eub-dject} s'opposer \tabularnewline
\en{produce} & \pron{pro-diouss} des produits & \pron{preu-diouss} produire \tabularnewline
\en{conduct} & \pron{kon-deukt} le comportement & \pron{keun-deukt} diriger \tabularnewline
\en{permit} & \pron{peu-mitt} un permis & \pron{peu-mitt} permettre \tabularnewline
\en{increase} & \pron{in-kriiss} une augmentation & \pron{in-kriiss} augmenter \tabularnewline
\end{longtable}
\end{tableau}

\section{60 paires minimales : l'entraînement de l'oreille}
Une \emph{paire minimale} est un couple de mots qui ne diffèrent que par un son. C'est
l'exercice fondamental pour installer les sons anglais dans votre oreille.

\textbf{Protocole :} faites écouter les deux mots par un outil audio (Cambridge, Forvo),
puis répétez-les 5 fois chacun, puis demandez à quelqu'un (ou à une application) de vous en
dire un seul au hasard : vous devez identifier lequel. 5 paires par jour suffisent.

\begin{tableau}[60 paires minimales --- par son]
\footnotesize
\begin{longtable}{@{}p{2.6cm}p{4.0cm}p{4.0cm}p{3.6cm}@{}}
\entete{Contraste} & \entete{Mot 1} & \entete{Mot 2} & \entete{Sens} \tabularnewline
\endfirsthead
\entete{Contraste} & \entete{Mot 1} & \entete{Mot 2} & \entete{Sens} \tabularnewline
\endhead
\hline
\ipaiot\ / \ipaee\ipalong & \en{ship} & \en{sheep} & bateau / mouton \tabularnewline
\ipaiot\ / \ipaee\ipalong & \en{sit} & \en{seat} & s'asseoir / siège \tabularnewline
\ipaiot\ / \ipaee\ipalong & \en{live} & \en{leave} & habiter / quitter \tabularnewline
\ipaiot\ / \ipaee\ipalong & \en{fill} & \en{feel} & remplir / sentir \tabularnewline
\ipaiot\ / \ipaee\ipalong & \en{slip} & \en{sleep} & glisser / dormir \tabularnewline
\ipaiot\ / \ipaee\ipalong & \en{bit} & \en{beat} & un morceau / battre \tabularnewline
\ipaiot\ / \ipaee\ipalong & \en{hit} & \en{heat} & frapper / chaleur \tabularnewline
\ipaiot\ / \ipaee\ipalong & \en{fit} & \en{feet} & convenir / pieds \tabularnewline
\hline
\ipae\ / \ipaE & \en{cat} & \en{cut} & chat / coupure \tabularnewline
\ipae\ / \ipaE & \en{bad} & \en{bed} & mauvais / lit \tabularnewline
\ipae\ / \ipaE & \en{man} & \en{men} & homme / hommes \tabularnewline
\ipae\ / \ipaE & \en{hat} & \en{hot} & chapeau / chaud \tabularnewline
\ipae\ / \ipaE & \en{pan} & \en{pen} & poêle / stylo \tabularnewline
\ipae\ / \ipaE & \en{bag} & \en{beg} & sac / mendier \tabularnewline
\ipae\ / \ipaE & \en{sat} & \en{set} & assis / un ensemble \tabularnewline
\hline
\ipalv\ / \ipaups & \en{cup} & \en{coop} & tasse / poulailler \tabularnewline
\ipalv\ / \ipaoo\ipalong & \en{full} & \en{fool} & plein / fou \tabularnewline
\ipalv\ / \ipaa\ipalong & \en{cut} & \en{cart} & couper / chariot \tabularnewline
\ipalv\ / \ipaer\ipalong & \en{fun} & \en{fern} & amusement / fougère \tabularnewline
\hline
\ipatheta\ / s & \en{think} & \en{sink} & penser / évier \tabularnewline
\ipatheta\ / s & \en{thick} & \en{sick} & épais / malade \tabularnewline
\ipatheta\ / f & \en{three} & \en{free} & trois / libre \tabularnewline
\ipatheta\ / t & \en{thing} & \en{ting} & chose / tintement \tabularnewline
\ipadh\ / d & \en{they} & \en{day} & ils / jour \tabularnewline
\ipadh\ / d & \en{then} & \en{den} & alors / tanière \tabularnewline
\ipadh\ / z & \en{this} & \en{zoo} & ceci / zoo \tabularnewline
\ipadh\ / d & \en{those} & \en{doze} & ceux-là / somnoler \tabularnewline
\hline
/w/ / v & \en{west} & \en{vest} & ouest / gilet \tabularnewline
/w/ / v & \en{wine} & \en{vine} & vin / vigne \tabularnewline
/w/ / v & \en{wale} & \en{veil} & --- / voile \tabularnewline
/w/ / v & \en{wet} & \en{vet} & mouillé / vétérinaire \tabularnewline
\hline
/b/ / p & \en{bat} & \en{pat} & chauve-souris / tape \tabularnewline
/b/ / p & \en{berry} & \en{perry} & baie / --- \tabularnewline
/b/ / p & \en{ban} & \en{pan} & interdiction / poêle \tabularnewline
/b/ / p & \en{boat} & \en{coat} & bateau / manteau \tabularnewline
/b/ / p & \en{back} & \en{pack} & dos / paquet \tabularnewline
\hline
/d/ / t & \en{doe} & \en{toe} & biche / orteil \tabularnewline
/d/ / t & \en{dear} & \en{tear} & cher / larme \tabularnewline
/d/ / t & \en{dime} & \en{time} & dix cents / temps \tabularnewline
\hline
/g/ / k & \en{goat} & \en{coat} & chèvre / manteau \tabularnewline
/g/ / k & \en{gift} & \en{kift} & cadeau / --- \tabularnewline
/g/ / k & \en{gale} & \en{kale} & grand vent / chou \tabularnewline
\hline
/n/ /\ipang/ & \en{sin} & \en{sing} & péché / chanter \tabularnewline
/n/ /\ipang/ & \en{thin} & \en{thing} & mince / chose \tabularnewline
/n/ /\ipang/ & \en{ran} & \en{rang} & courut / sonna \tabularnewline
/n/ /\ipang/ & \en{ban} & \en{bang} & interdiction / bang \tabularnewline
\hline
/l/ / r & \en{light} & \en{right} & léger / droite \tabularnewline
/l/ / r & \en{lead} & \en{read} & mener / lire \tabularnewline
/l/ / r & \en{love} & \en{rove} & amour / errer \tabularnewline
/l/ / r & \en{fly} & \en{fry} & voler / frire \tabularnewline
/l/ / r & \en{collect} & \en{correct} & collecter / corriger \tabularnewline
/l/ / r & \en{low} & \en{row} & bas / ramer \tabularnewline
\hline
/z/ / s & \en{zip} & \en{sip} & fermeture éclair / gorgée \tabularnewline
/z/ / s & \en{bus} & \en{buzz} & bus / bourdonner \tabularnewline
/z/ / s & \en{rice} & \en{rise} & riz / montée \tabularnewline
/z/ / s & \en{place} & \en{plays} & lieu / joue \tabularnewline
\hline
/\ipash/ / s & \en{she} & \en{sea} & elle / mer \tabularnewline
/\ipash/ / t\ipash/ & \en{share} & \en{chair} & partager / chaise \tabularnewline
/\ipash/ / t\ipash/ & \en{shore} & \en{chore} & rivage / corvée \tabularnewline
/\ipazh/ / d\ipazh/ & \en{vision} & \en{Vishnu} & vision / --- \tabularnewline
\hline
\ipaou / \ipao\ipalong & \en{goat} & \en{got} & chèvre / obtenu \tabularnewline
\ipaou / \ipao & \en{no} & \en{now} & non / maintenant \tabularnewline
\ipaou / \ipaoo\ipalong & \en{hope} & \en{hoop} & espoir / cerceau \tabularnewline
/\ipaa\ipalong/ / \ipae/ & \en{cart} & \en{cat} & chariot / chat \tabularnewline
/\ipao\ipalong/ / \ipao\ (bref)/ & \en{caught} & \en{cot} & attrapé / lit d'enfant \tabularnewline
/\ipaiot\ipeschwa/ / \ipaE\ipeschwa/ & \en{here} & \en{hair} & ici / cheveux \tabularnewline
/\ipaiot\ipeschwa/ / \ipaE\ipeschwa/ & \en{near} & \en{nair} & près / --- \tabularnewline
/h/ / voyelle & \en{hill} & \en{ill} & colline / malade \tabularnewline
/h/ / voyelle & \en{heat} & \en{eat} & chaleur / manger \tabularnewline
/h/ / voyelle & \en{hat} & \en{at} & chapeau / à \tabularnewline
\end{longtable}
\end{tableau}

\section{Exercices}

\begin{exo}[5.1 --- Identifiez le son]
Écoutez (ou imaginez la prononciation correcte) et classez chaque mot dans la bonne colonne :
\en{ship -- sheep -- live -- leave -- feel -- fill -- sit -- seat -- heat -- hit -- rich -- reach}
\begin{center}
\begin{tabular}{@{}p{5.5cm}p{5.5cm}@{}}
\entete{/\ipaiot/ bref} & \entete{/\ipaee\ipalong/ long} \tabularnewline
\hline
\trou[4.5cm] & \trou[4.5cm] \\
\trou[4.5cm] & \trou[4.5cm] \\
\end{tabular}
\end{center}
\end{exo}
\begin{corrige}[5.1]
\textbf{Bref} /\ipaiot/ : \en{ship, live, fill, sit, hit, rich}.
\textbf{Long} /\ipaee\ipalong/ : \en{sheep, leave, feel, seat, heat, reach}.
\end{corrige}

\begin{exo}[5.2 --- Transcrivez à la française]
Donnez la prononciation approchée des douze mots suivants :
\en{water, thought, three, this, house, know, right, work, book, food, cup, go}.
\end{exo}
\begin{corrige}[5.2]
\en{water} \pron{ouo-teu} \quad \en{thought} \pron{thoot} \quad \en{three} \pron{thrii} \quad
\en{this} \pron{dhis} \quad \en{house} \pron{ha-ouss} \quad \en{know} \pron{nou} \quad
\en{right} \pron{raït} \quad \en{work} \pron{oueu-k} \quad \en{book} \pron{bouk} \quad
\en{food} \pron{fouud} \quad \en{cup} \pron{keup} \quad \en{go} \pron{gou-ou}.
\end{corrige}

\begin{exo}[5.3 --- Quelle paire ?]
Pour chaque phrase, choisissez le mot correct :
\begin{enumerate}
  \item I want to (sit / seat) down.
  \item The (ship / sheep) is very big on the sea.
  \item I (think / sink) you are right.
  \item This is my (wine / vine) --- it's from France.
  \item She has beautiful (hair / here).
  \item Don't (fill / feel) bad about it.
  \item The (cart / cut) is in the garden.
  \item We need a (bag / beg) for the shopping.
  \item He will (live / leave) tomorrow morning.
  \item My (dad / dead) works in a bank.
\end{enumerate}
\end{exo}
\begin{corrige}[5.3]
1. sit \quad 2. ship \quad 3. think \quad 4. wine \quad 5. hair \quad
6. feel \quad 7. cart \quad 8. bag \quad 9. leave \quad 10. dad
\end{corrige}

\begin{exo}[5.4 --- Les lettres muettes]
Dans chaque mot, entourez la lettre qui ne se prononce pas :
\en{knee, write, climb, listen, walk, hour, sign, guess, half, know, comb, island, talk, honest, wrist}.
\end{exo}
\begin{corrige}[5.4]
\en{knee (k), write (w), climb (b), listen (t), walk (l), hour (h), sign (g), guess (u),
half (l), know (k), comb (b), island (s), talk (l), honest (h), wrist (w)}.
\end{corrige}

\begin{mission}[Chapitre 5 --- 10 jours de sons]
\textbf{Jour 1-2 :} les voyelles longues et brèves (paires \en{sheep/ship, full/fool, cart/cut}).
\textbf{Jour 3-4 :} \ipatheta\ et \ipadh\ (30 mots, 5 fois chacun devant un miroir, langue visible entre les dents).
\textbf{Jour 5 :} /h/ aspiré (30 mots) --- expirez avant de parler.
\textbf{Jour 6 :} /w/ contre /v/ (20 paires).
\textbf{Jour 7 :} /r/ anglais (30 mots devant un miroir : la langue ne touche rien).
\textbf{Jour 8 :} \ipang\ final (\en{sing, long, thing, bring, wrong, song, young, strong}).
\textbf{Jour 9 :} les lettres muettes (tableau ci-dessus, 60 mots).
\textbf{Jour 10 :} test --- lisez à voix haute les 60 paires minimales, enregistrez-vous,
comparez avec un audio de référence, notez vos 3 sons les plus faibles. Ce sont vos
priorités du chapitre 7.
\end{mission}


% =====================================================================
\chapter{Accent tonique, rythme, intonation et \en{connected speech}}

Ce chapitre explique \textbf{pourquoi les francophones ne comprennent pas les anglophones
natifs} --- et ce n'est pas une question de vocabulaire.

\section{Deux rythmes incompatibles}
\begin{tableau}[Français contre anglais : la différence fondamentale]
\begin{tabularx}{\linewidth}{@{}p{4.0cm}X@{}}
\entete{} & \entete{Caractéristique} \tabularnewline
\hline
\textbf{Français} & Langue à \textbf{rythme syllabique} : chaque syllabe a à peu près la
même durée et la même force. \en{Je-ne-sais-pas} = quatre coups réguliers. L'accent tombe
sur la dernière syllabe du groupe. \tabularnewline
\hline
\textbf{Anglais} & Langue à \textbf{accent tonique} : les syllabes accentuées sont
\textbf{fortes, longues et hautes}, les autres sont \textbf{écrasées, courtes et réduites}
au schwa. Le temps entre deux accents est à peu près constant, quel que soit le nombre de
syllabes faibles entre elles. \tabularnewline
\hline
\textbf{Conséquence} & Un anglophone qui entend du français régulier le trouve « plat,
rapide, monotone ». Un francophone qui entend de l'anglais accentué trouve que les mots
« disparaissent » : il cherche \en{I want to go to the shop} et entend \en{I wanna go t'the
shop}. \textbf{Vous ne comprenez pas parce que vous cherchez des syllabes françaises là où
il n'y a que des accents anglais.} \tabularnewline
\end{tabularx}
\end{tableau}

\begin{memo}[La règle d'or du chapitre]
En anglais, ce qui compte n'est pas ce que vous dites, c'est \textbf{où vous frappez}.
Une syllabe accentuée mal placée rend un mot incompréhensible, même si tous ses sons sont
justes. Exemple : \en{banana} \pron{be-NA-neu} est compris ; \pron{ba-na-na} (à la
française) ne l'est pas.
\end{memo}

\section{L'accent tonique dans les mots}

\subsection{Les quatre principes}
\begin{enumerate}
  \item Tout mot anglais de deux syllabes ou plus a \textbf{une} syllabe accentuée
  (éventuellement une secondaire).
  \item Une syllabe accentuée est \textbf{plus longue, plus forte et plus aiguë}.
  \item Toute syllabe non accentuée tend vers le \textbf{schwa} /\ipeschwa/.
  \item L'accent est \textbf{imprévisible} --- il faut donc l'apprendre avec le mot, comme
  une partie du mot. Dans votre Anki, notez toujours la syllabe accentuée (par exemple en
  MAJUSCULES : \en{baNAna}, \en{COMfortable}, \en{deVELopment}).
\end{enumerate}

\subsection{Les règles qui marchent 80\,\% du temps}
\begin{tableau}[Règles de placement de l'accent]
\begin{tabularx}{\linewidth}{@{}p{5.6cm}X@{}}
\entete{Règle} & \entete{Exemples} \tabularnewline
\hline
\textbf{Nom de 2 syllabes} : accent sur la 1\textsuperscript{re} &
\en{TAble, WAter, MOney, CIty, PREsent, EXport, CHIna} \tabularnewline
\textbf{Verbe de 2 syllabes} : accent sur la 2\textsuperscript{e} &
\en{to beGIN, to preSENT, to exPORT, to deCIDE, to reLAX, to rePEAT} \tabularnewline
\textbf{Mots en \en{-tion}, \en{-sion}, \en{-cian}} : accent sur l'avant-dernière syllabe &
\en{naTION, inforMAtion, deciSIon, musIcian, televiSIon} \tabularnewline
\textbf{Mots en \en{-ic}} : accent sur l'avant-dernière &
\en{scienTIfic, draMAtic, econoMIc, fanTAStic, reaLIStic} \tabularnewline
\textbf{Mots en \en{-ity}, \en{-ify}, \en{-ise}} : accent 3 syllabes avant la fin &
\en{uniVERsity, possiBIlity, CLAssify, adVERtise, ORganise} \tabularnewline
\textbf{Mots en \en{-ee}, \en{-eer}, \en{-esque}} : accent \textbf{sur} la finale &
\en{emPLOyee, volunTEER, engiNEER, pictuRESque} \tabularnewline
\textbf{Composés noms} : accent sur le 1\textsuperscript{er} élément &
\en{GREENhouse, FOOTball, POSTman, BATHroom} \tabularnewline
\textbf{Nombres en \en{-teen}} : accent sur la 2\textsuperscript{e} ; \en{-ty} : sur la 1\textsuperscript{re} &
\en{thirTEEN} contre \en{THIRty} --- \textbf{différence capitale à l'oral} \tabularnewline
\end{tabularx}
\end{tableau}

\begin{piege}[\en{thirteen} ou \en{thirty} ?]
C'est la confusion la plus coûteuse de la vie quotidienne (prix, dates, numéros, âges).
\en{thirteen} = \pron{theu-tiin} (accent final, \en{ee} long) ; \en{thirty} =
\pron{THEU-di} (accent initial, \en{ir} bref). En cas de doute, l'anglophone demande :
\en{``Sorry, thirteen or thirty?''} --- entraînez la paire 13/30, 14/40, 15/50, 16/60,
17/70, 18/80, 19/90 jusqu'à l'automatisme.
\end{piege}

\subsection{Accent de mot et sens}
Certains mots changent de \textbf{sens} selon l'accent (voir tableau des \en{heteronyms},
chapitre 5). Retenez les dix plus utiles : \en{record, present, object, produce, conduct,
permit, increase, decrease, protest, suspect}.

\section{Le schwa : le son qui change tout}
/\ipeschwa/ est une voyelle neutre, à peine articulée, qui remplace \textbf{toute voyelle
d'une syllabe non accentuée}. C'est le son le plus fréquent de l'anglais.

\begin{tableau}[Le schwa se cache partout --- mêmes lettres, son réduit]
\footnotesize
\begin{longtable}{@{}p{3.4cm}p{3.0cm}p{8.0cm}@{}}
\entete{Mot} & \entete{Prononciation} & \entete{Où sont les schwas ?} \tabularnewline
\endfirsthead
\entete{Mot} & \entete{Prononciation} & \entete{Où sont les schwas ?} \tabularnewline
\endhead
\hline
\en{about} & \pron{e-baout} & Le \en{a} initial devient schwa. \tabularnewline
\en{banana} & \pron{be-NA-neu} & Les deux \en{a} non accentués deviennent schwa. \tabularnewline
\en{problem} & \pron{PRO-beulm} & Le \en{e} de la 2\textsuperscript{e} syllabe. \tabularnewline
\en{support} & \pron{se-POOT} & Le \en{u} initial. \tabularnewline
\en{photograph} & \pron{FOU-teu-graaf} & Le \en{o} de la 2\textsuperscript{e} syllabe. \tabularnewline
\en{photographer} & \pron{fe-TO-gra-feu} & Les trois voyelles non accentuées. \tabularnewline
\en{computer} & \pron{keum-Piou-teu} & \en{o} et \en{e} finaux. \tabularnewline
\en{interesting} & \pron{IN-tres-ting} & Souvent réduit à 3 syllabes, pas 5. \tabularnewline
\en{comfortable} & \pron{KOMF-teu-bl} & 3 syllabes, pas 5. \tabularnewline
\en{vegetable} & \pron{VEDJ-teu-bl} & 3 syllabes, pas 4. \tabularnewline
\en{today} & \pron{t'e-dé-i} & Le \en{o} devient \pron{t'}. \tabularnewline
\en{police} & \pron{pe-Liiss} & Le \en{o} initial. \tabularnewline
\end{longtable}
\end{tableau}

\begin{astuce}[Exercice du schwa --- 5 min/jour]
Prenez n'importe quel mot de trois syllabes ou plus, et dites-le de deux façons :
\textbf{(1)} à la française, toutes les syllabes égales ; \textbf{(2)} à l'anglaise, en
frappant la syllabe accentuée et en écrasant les autres. Enregistrez les deux versions.
Vous entendrez immédiatement la différence --- et vous comprendrez pourquoi les natifs ne
vous comprenaient pas.
\end{astuce}

\section{L'accent de phrase : les mots forts et les mots faibles}
Dans une phrase anglaise, on accentue les \textbf{mots de contenu} et on écrase les
\textbf{mots de forme}.

\begin{tableau}[Mots accentués et mots réduits]
\begin{tabularx}{\linewidth}{@{}p{7.5cm}X@{}}
\entete{\textbf{Accentués} (portent le sens)} & \entete{\textbf{Réduits} (structure)} \tabularnewline
\hline
Noms, verbes pleins, adjectifs, adverbes ; pronoms démonstratifs (\en{this, that}) ;
négations (\en{not, never}) ; mots interrogatifs (\en{what, where, why, how}) ;
auxiliaires en fin de phrase ou contractés avec emphase ; prépositions suivies d'un accent
de contraste. &
Articles (\en{a, an, the}) ; prépositions (\en{of, to, for, at, from, by}) ; conjonctions
(\en{and, but, or, as, that}) ; auxiliaires (\en{be, have, do, can, will, would}) ;
pronoms personnels (\en{he, him, her, them, us, you}) ; possessifs (\en{my, your, his}). \tabularnewline
\end{tabularx}
\end{tableau}

\ex{I want to go to the shop to buy some bread.}{Je veux aller à la boutique acheter du pain.}
\textbf{Version réelle :} \pron{aï WANT t' g' t' the SHOP t' BUY s'm BRED} ---
cinq accents forts, tout le reste écrasé.

\ex{She can speak French and English.}{Elle sait parler français et anglais.}
\textbf{Version réelle :} \pron{ch' k'n SPIIK French 'n INGLISH}.

\begin{piege}[La négation et le contraste doivent être frappés fort]
\en{I \textbf{didn't} say that} : si vous réduisez \en{didn't}, on entend \en{I said that}
--- l'inverse du sens. De même : \en{I can do it} (\en{can} faible \pron{k'n}) contre
\en{I \textbf{can't} do it} (\en{can't} fort \pron{kaant}). Entraînez-vous à frapper
\en{not, never, no, can't, don't, isn't, wasn't}.
\end{piege}

\section{L'intonation : la mélodie des phrases}
\begin{tableau}[Les 5 courbes d'intonation essentielles]
\begin{tabularx}{\linewidth}{@{}p{4.2cm}p{2.6cm}X@{}}
\entete{Type de phrase} & \entete{Mélodie} & \entete{Exemple} \tabularnewline
\hline
Question fermée (oui/non) & Montante \ra & \en{Are you coming?} \pron{a y' CO-ming $\monte$} \tabularnewline
Question ouverte (\en{wh-}) & \textbf{Descendante} \descend & \en{Where do you live?} --- \textbf{erreur fréquente : les francophones la font monter.} \tabularnewline
Affirmation & Descendante \descend & \en{I live in Dakar.} \pron{i LIVE in DA-kaar $\descend$} \tabularnewline
Liste & Montante, montante\ldots\ puis descendante & \en{I bought bread$\monte$, milk$\monte$ and eggs$\descend$.} \tabularnewline
Surprise, doute, politesse & Très montante & \en{Really?} \pron{rii-AL-lii $\monte$} \quad \en{Sorry?} \tabularnewline
\end{tabularx}
\end{tableau}

\begin{astuce}[L'intonation porte le message émotionnel]
La même phrase dite avec trois mélodies donne trois sens :\\
\en{``You're coming.''} \descend\ = constatation. \quad
\en{``You're coming?''} \monte\ = question étonnée. \quad
\en{``\textbf{You're} coming?''} \monte\ avec accent sur \en{you} = « \emph{toi}, tu viens ? ».
\end{astuce}

\section{\en{Connected speech} : pourquoi les mots disparaissent}
Les anglophones ne disent pas les mots un par un : ils les enchaînent en un seul flux sonore.
Voici les sept mécanismes --- chacun avec un exemple que vous devez savoir reconnaître à l'oreille.

\subsection{1. Le \emph{linking} (consonne finale + voyelle initiale)}
La consonne finale d'un mot colle à la voyelle du mot suivant, comme une liaison française.
\begin{itemize}
  \item \en{an apple} \pron{e-NA-pl} --- le \en{n} passe dans \en{apple}.
  \item \en{turn it off} \pron{teu-ri-NI-tof}.
  \item \en{not at all} \pron{no-ta-TAAL}.
  \item \en{pick it up} \pron{pi-ki-TUP}.
  \item \en{an hour ago} \pron{e-nau-e-ROU}.
\end{itemize}

\subsection{2. Le \en{r} de liaison (britannique)}
Le \en{r} final, muet en britannique, se prononce devant une voyelle :
\en{far away} \pron{faa-re-wé-i}, \en{for example} \pron{fo-rig-Zaampl},
\en{here it is} \pron{hi-e-ri-TIZ}.

\subsection{3. Les intrusions /w/ et /j/}
Entre deux voyelles, l'anglais insère un son de transition :
\begin{itemize}
  \item \en{go out} \pron{gou-Waout} ; \en{do it} \pron{dou-WIT} ; \en{who is it?} \pron{hou-WI-zit}.
  \item \en{I am} \pron{aï-YAM} ; \en{see it} \pron{sii-YIT} ; \en{the end} \pron{dhi-YEND}.
\end{itemize}

\subsection{4. Les élisions (sons supprimés)}
\begin{itemize}
  \item \en{next day} \pron{nèks-dé-i} (le \en{t} saute).
  \item \en{last night} \pron{laas-naït}.
  \item \en{handbag} \pron{hann-bag}.
  \item \en{comfortable} \pron{KOMF-teu-bl} ; \en{family} \pron{FAM-li} ; \en{camera} \pron{KAM-reu}.
  \item \en{probably} \pron{PROB-li} ; \en{actually} \pron{AK-chou-a-li} ; \en{chocolate} \pron{TCHOK-l't}.
\end{itemize}

\subsection{5. Les assimilations (sons qui se transforment)}
\begin{itemize}
  \item \en{t} + \en{y} \ra \en{ch} : \en{don't you} \pron{doun-tchou} ; \en{meet you} \pron{mii-tchou}.
  \item \en{d} + \en{y} \ra \en{j} : \en{did you} \pron{di-djou} ; \en{would you} \pron{wou-djou}.
  \item \en{s} + \en{y} \ra \en{ch} : \en{miss you} \pron{mi-chou}.
  \item \en{z} + \en{y} \ra \en{j} : \en{was your} \pron{ouo-djou}.
\end{itemize}

\subsection{6. Le \emph{glottal stop} (coup de glotte)}
Le \en{t} entre voyelles ou devant consonne devient un arrêt net de la gorge (très britannique) :
\en{butter} \pron{ba-eu}, \en{football} \pron{fou-bo-l}, \en{not now} \pron{no-naou},
\en{water} \pron{ouo-eu}. En américain, ce \en{t} devient plutôt un \en{d} bref
(\en{water} = \pron{oua-deu}).

\subsection{7. Les formes faibles (\emph{weak forms})}
\begin{tableau}[Formes faibles des petits mots --- à reconnaître absolument]
\footnotesize
\begin{longtable}{@{}p{2.4cm}p{3.2cm}p{2.6cm}p{5.8cm}@{}}
\entete{Mot} & \entete{Forme forte} & \entete{Forme faible} & \entete{Exemple} \tabularnewline
\endfirsthead
\entete{Mot} & \entete{Forte} & \entete{Faible} & \entete{Exemple} \tabularnewline
\endhead
\hline
\en{and} & \pron{and} & \pron{'n} & \en{fish and chips} = \pron{fi-ch'n tchips} \tabularnewline
\en{of} & \pron{ov} & \pron{euv} (schwa) & \en{a cup of tea} = \pron{e-keup-euv-tii} \tabularnewline
\en{to} & \pron{tou} & \pron{t'} & \en{go to bed} = \pron{gou-t'-bèd} \tabularnewline
\en{for} & \pron{foo} & \pron{fe} & \en{for me} = \pron{fe-mii} \tabularnewline
\en{can} & \pron{kann} & \pron{k'n} & \en{I can swim} = \pron{aï k'n SWIM} \tabularnewline
\en{have} & \pron{hav} & \pron{hev/'v} & \en{should have} = \pron{chou-dev} \ra \en{shoulda} \tabularnewline
\en{do} & \pron{dou} & \pron{de} & \en{do you know?} = \pron{d'ye NOU} \tabularnewline
\en{at} & \pron{at} & \pron{et} & \en{at work} = \pron{et-oueu-k} \tabularnewline
\en{from} & \pron{from} & \pron{freum} & \en{from Dakar} = \pron{freum-da-kaar} \tabularnewline
\en{them} & \pron{dhèm} & \pron{dhem/'em} & \en{tell them} = \pron{tè-l'em} \tabularnewline
\en{us} & \pron{euss} & \pron{es} & \en{let us go} = \en{let's go} = \pron{lèts-gou} \tabularnewline
\en{was} & \pron{ouoz} & \pron{weu-z} & \en{it was good} = \pron{i-teu-z-GOUD} \tabularnewline
\en{your} & \pron{yoo} & \pron{yeu} & \en{your name} = \pron{yeu-né-im} \tabularnewline
\en{he/him/her} & \pron{hii/im/heu} & \pron{i/im/eu} & \en{tell him} = \pron{té-lim} ; \en{I saw her} = \pron{aï soo-weu} \tabularnewline
\end{longtable}
\end{tableau}

\subsection{Les réductions familières (à comprendre, à utiliser avec prudence)}
\begin{tableau}[Réductions courantes de l'anglais parlé]
\footnotesize
\begin{longtable}{@{}p{3.2cm}p{3.0cm}p{8.0cm}@{}}
\entete{Forme réduite} & \entete{Forme complète} & \entete{Exemple et traduction} \tabularnewline
\endfirsthead
\entete{Réduite} & \entete{Complète} & \entete{Exemple} \tabularnewline
\endhead
\hline
\en{gonna} & \en{going to} & \en{I'm gonna call you.} --- \textit{Je vais t'appeler.} \tabularnewline
\en{wanna} & \en{want to} & \en{Do you wanna eat?} --- \textit{Tu veux manger ?} \tabularnewline
\en{gotta} & \en{got to / have to} & \en{I gotta go.} --- \textit{Je dois y aller.} \tabularnewline
\en{hafta} & \en{have to} & \en{We hafta leave now.} --- \textit{On doit partir maintenant.} \tabularnewline
\en{oughta} & \en{ought to} & \en{You oughta see a doctor.} --- \textit{Tu devrais voir un médecin.} \tabularnewline
\en{supposta} & \en{supposed to} & \en{What was I supposta do?} --- \textit{Qu'est-ce que j'étais censé faire ?} \tabularnewline
\en{used to} & \en{used to} & \pron{iou-s-tou} --- \textit{avant, j'avais l'habitude de\ldots} \tabularnewline
\en{lemme} & \en{let me} & \en{Lemme see.} --- \textit{Laisse-moi voir.} \tabularnewline
\en{gimme} & \en{give me} & \en{Gimme a minute.} --- \textit{Donne-moi une minute.} \tabularnewline
\en{dunno} & \en{don't know} & \en{I dunno.} --- \textit{Je sais pas.} \tabularnewline
\en{kinda / sorta} & \en{kind of / sort of} & \en{It's kinda cold.} --- \textit{C'est un peu froid.} \tabularnewline
\en{d'you / dya} & \en{do you} & \en{D'you want some?} --- \textit{Tu en veux ?} \tabularnewline
\en{whaddaya} & \en{what do you} & \en{Whaddaya mean?} --- \textit{Qu'est-ce que tu veux dire ?} \tabularnewline
\en{c'mon} & \en{come on} & \en{C'mon, let's go!} --- \textit{Allez, on y va !} \tabularnewline
\en{y'know} & \en{you know} & \en{It's hard, y'know.} --- \textit{C'est dur, tu sais.} \tabularnewline
\en{innit} \en{(BrE)} & \en{isn't it} & \en{It's cold, innit?} --- \textit{Il fait froid, non ?} \tabularnewline
\end{longtable}
\end{tableau}

\begin{astuce}[Les comprendre, oui ; les utiliser, plus tard]
Vous devez \textbf{reconnaître} toutes ces réductions à l'oreille dès le mois 3 --- sinon les
séries resteront incompréhensibles. Mais ne les \textbf{employez} qu'à partir du mois 8 ou 9,
et seulement dans un contexte familier : un \en{gonna} dans un entretien d'embauche est
maladroit, dans une discussion entre amis il est naturel.
\end{astuce}

\section{Exercices de rythme et d'intonation}

\begin{exo}[6.1 --- Marquez les accents]
Écrivez les mots suivants en majuscules sur la syllabe accentuée :
\en{banana, information, comfortable, development, photographer, university, employee,
interesting, television, tomorrow, yesterday, computer}.
\end{exo}
\begin{corrige}[6.1]
\en{baNAna, inforMAtion, KOMfortable, deVElopment, phoTOgrapher, uniVERsity, employEE,
INteresting, TELevision, toMOrrow, YEStarday, comPUter}.
\end{corrige}

\begin{exo}[6.2 --- Nom ou verbe ?]
Prononcez correctement selon le sens, puis traduisez :
\begin{enumerate}
  \item I bought a new \en{record}. / Please \en{record} my voice.
  \item She gave me a \en{present}. / He will \en{present} the project.
  \item There was an \en{increase} in prices. / Prices \en{increase} every year.
  \item Do you \en{object}? / It's a strange \en{object}.
\end{enumerate}
\end{exo}
\begin{corrige}[6.2]
1. \pron{RÈ-kood} un disque / \pron{ri-KOOD} enregistrer. 2. \pron{PRÈ-z'nt} un cadeau /
\pron{pri-ZÈNT} présenter. 3. \pron{IN-kriiss} une augmentation / \pron{in-KRISS} augmenter.
4. \pron{eub-DJECT} s'opposer / \pron{OB-dject} un objet.
\end{corrige}

\begin{exo}[6.3 --- Écrivez ce que vous entendez]
Remettez ces phrases réduites en anglais correct et complet :
\begin{enumerate}
  \item \en{``I'm gonna call ya later.''}
  \item \en{``D'you wanna cuppa tea?''}
  \item \en{``He shoulda told us.''}
  \item \en{``Whaddaya doin'?''}
  \item \en{``Lemme see it, gimme a sec.''}
  \item \en{``I dunno what he's talkin' about.''}
  \item \en{``We gotta go, it's kinda late.''}
  \item \en{``Tell 'em I said hi.''}
\end{enumerate}
\end{exo}
\begin{corrige}[6.3]
\begin{enumerate}
  \item \en{I am going to call you later.} --- \textit{Je t'appellerai plus tard.}
  \item \en{Do you want a cup of tea?} --- \textit{Tu veux une tasse de thé ?}
  \item \en{He should have told us.} --- \textit{Il aurait dû nous le dire.}
  \item \en{What are you doing?} --- \textit{Qu'est-ce que tu fais ?}
  \item \en{Let me see it, give me a second.} --- \textit{Laisse-moi voir, donne-moi une seconde.}
  \item \en{I don't know what he is talking about.} --- \textit{Je ne sais pas de quoi il parle.}
  \item \en{We have got to go, it is kind of late.} --- \textit{On doit y aller, il est un peu tard.}
  \item \en{Tell them I said hello.} --- \textit{Dis-leur que je leur ai dit bonjour.}
\end{enumerate}
\end{corrige}

\begin{exo}[6.4 --- Intonation]
Indiquez la mélodie (montante \monte\ ou descendante \descend) :
\begin{enumerate}
  \item \en{Do you speak English?}
  \item \en{Where do you come from?}
  \item \en{My name is Ali.}
  \item \en{I need bread, milk and eggs.}
  \item \en{You're joking!}
  \item \en{Can you help me?}
  \item \en{What time is it?}
  \item \en{It's half past nine.}
\end{enumerate}
\end{exo}
\begin{corrige}[6.4]
1. \monte\ \ 2. \descend\ \ 3. \descend\ \ 4. \monte\monte\descend\ (sur \en{eggs}) \
5. \monte\ \ 6. \monte\ \ 7. \descend\ \ 8. \descend
\end{corrige}

\begin{mission}[Chapitre 6 --- le drill des 5 accents]
Pendant 7 jours, prenez chaque jour 5 phrases de ce livre (n'importe lesquelles) et :
\textbf{(1)} marquez au crayon les syllabes accentuées ; \textbf{(2)} lisez-les en exagérant
les accents (fort et haut) et en écrasant le reste ; \textbf{(3)} enregistrez-vous ;
\textbf{(4)} comparez à un natif (YouGlish : tapez la phrase, écoutez 5 personnes la dire).
À la fin de la semaine, votre rythme aura changé --- les gens vous diront que vous
« sonnez mieux », sans savoir pourquoi.
\end{mission}


% =====================================================================
\chapter{Clinique de prononciation pour francophones}

Voici les quinze erreurs systématiques des francophones (dont les Africains francophones),
avec pour chacune : la cause, la correction et un entraînement. C'est le chapitre le plus
rentable du livre : chacune de ces erreurs corrigée améliore immédiatement votre
intelligibilité.

\section{Les 15 erreurs et leur correction}

\subsection{Erreur 1 --- Avaler le \en{h}}
\textbf{Ce que vous dites :} \en{he} \pron{i}, \en{house} \pron{ouss}, \en{happy} \pron{api}.
\textbf{Ce qu'on entend :} \en{eat, ounce, appy} --- ou simplement « un accent étrange ».
\textbf{Cause :} le \en{h} français est muet, votre bouche a pris l'habitude de ne pas souffler.
\textbf{Correction :} le \en{h} anglais est un vrai souffle. Expirez avant de prononcer la
voyelle, comme pour faire de la buée sur une vitre.
\textbf{Drill :} \en{he, him, her, his, here, house, home, hope, hand, head, hair, hat, hot,
heart, help, hold, how, hello, behind, ahead, perhaps} --- 3 fois chacun, main devant la
bouche pour sentir le souffle.

\subsection{Erreur 2 --- Transformer le \en{th} en \en{s}, \en{z}, \en{f} ou \en{d}}
\textbf{Ce que vous dites :} \en{think} \pron{sink}, \en{this} \pron{dis}, \en{three}
\pron{free}, \en{brother} \pron{brozeu}.
\textbf{Correction :} langue \emph{entre} les dents --- vraiment entre, on doit la voir.
\ipatheta\ = souffle sans voix (\en{think}) ; \ipadh\ = souffle avec voix (\en{this}).
\textbf{Drill :} \en{think, thing, three, thank, thought, through, thousand, thumb, bath,
mouth, month, birthday} --- puis \en{this, that, these, those, there, they, them, then,
than, mother, father, brother, other, weather, with}.
\textbf{Astuce visuelle :} répétez devant un miroir jusqu'à voir votre langue dépasser.

\subsection{Erreur 3 --- Le \en{w} prononcé \en{v}}
\textbf{Ce que vous dites :} \en{we} \pron{vi}, \en{water} \pron{vateu}, \en{work} \pron{veuk}.
\textbf{Correction :} /w/ = lèvres très arrondies et projetées en avant, puis ouverture
rapide --- comme un « ou » qui glisse. Aucune lèvre sur les dents (c'est cela, le \en{v}).
\textbf{Drill :} \en{we, will, would, water, work, world, want, when, where, why, which,
woman, window, week, way, well, always, now, know, grow, show} --- en alternance avec
\en{ve, veal, very, vest, vet, vine, visit, voice}.

\subsection{Erreur 4 --- Le \en{r} de gorge français}
\textbf{Correction :} le \en{r} anglais ne vient pas de la gorge. Recroquevillez la langue
au milieu de la bouche, sans contact, lèvres légèrement arrondies. C'est un son « mou »,
proche d'un \en{w} très bref.
\textbf{Drill :} \en{red, right, run, rain, road, room, read, write, rice, ring, round,
river, around, sorry, really, very, every, remember, important, world, work, girl, bird}.
\textbf{Test :} dites \en{``around the world''} dix fois : si votre gorge gratte, recommencez.

\subsection{Erreur 5 --- Ajouter un \en{e} final (le fameux « euh »)}
\textbf{Ce que vous dites :} \en{book} \pron{bou-keu}, \en{dog} \pron{do-gueu}, \en{stop}
\pron{sto-peu}, \en{work} \pron{oueu-keu}.
\textbf{Cause :} le français n'aime pas finir sur une consonne ; votre bouche ajoute une
voyelle de soutien.
\textbf{Correction :} arrêtez le son net. Fermez la bouche à la fin de \en{book},
coupez le souffle à la fin de \en{stop}.
\textbf{Drill :} \en{book, look, take, make, work, walk, talk, stop, jump, help, want, need,
good, bad, big, small, cat, dog, map, cup, black, white, red, green, blue} --- en terminant
chaque mot par un silence franc.

\subsection{Erreur 6 --- Ne pas distinguer \en{ship} et \en{sheep}}
Voir chapitre 5 : bref contre long, tendu contre relâché. Drill : les 20 paires en
/\ipaiot/ -- /\ipaee\ipalong/ de la liste des paires minimales.

\subsection{Erreur 7 --- \en{no} prononcé \pron{no}}
\textbf{Correction :} le \en{o} anglais final est une diphtongue : \pron{nou-ou}.
Idem pour \en{go, so, home, phone, only, open, close, cold, both, note}.
\textbf{Drill :} \en{no, go, so, home, phone, hope, note, both, cold, only, open, close,
over, road, show, know, slow, grow, throw, alone}.

\subsection{Erreur 8 --- L'intonation plate et le rythme syllabique}
Vous parlez anglais avec le rythme du français : toutes les syllabes égales. Résultat : les
natifs doivent faire un effort pour vous comprendre, et vous fatiguez vite.
\textbf{Correction :} exagérez. Frappez les mots de contenu, écrasez les petits mots.
C'est inconfortable au début --- c'est exactement ce qu'il faut.
\textbf{Drill :} shadowing quotidien (chapitre 3) sur des phrases courtes.

\subsection{Erreur 9 --- Le \en{-s} final oublié ou mal sonorisé}
\textbf{Ce que vous dites :} \en{he work} au lieu de \en{he works}, \en{two book}.
\textbf{Et le pire :} le \en{-s} du pluriel se prononce /s/, /z/ ou /\ipaiot z/ selon le mot :
\en{cats} \pron{kats}, \en{dogs} \pron{dogz}, \en{boxes} \pron{bok-siz}.
\textbf{Correction :} après une consonne sourde \ra /s/ ; après une voyelle ou une consonne
sonore \ra /z/ ; après \en{s, z, sh, ch, x} \ra /iz/.
\textbf{Drill :} \en{books / bags / buses ; cats / dogs / houses ; maps / pens / watches}.

\subsection{Erreur 10 --- Le \en{ng} final prononcé \pron{ingue}}
\en{sing} = \pron{si-ng}, pas \pron{si-ngue}. La langue touche le palais, le son est nasal,
et \textbf{rien ne suit}. Drill : \en{sing, song, long, strong, bring, thing, wrong, young,
king, ring, morning, evening, nothing, everything, interesting, English}.

\subsection{Erreur 11 --- Les voyelles longues abrégées}
\en{car} = \pron{kaa}, pas \pron{kar}. \en{four} = \pron{foo-r}. \en{bird} = \pron{beuud}.
Le français n'a pas de voyelles longues ; l'anglais en tire des distinctions de sens.
Drill : tenez chaque voyelle une seconde : \en{car, park, heart, start, father} ---
\en{four, more, door, floor, for} --- \en{bird, word, world, girl, third}.

\subsection{Erreur 12 --- Le \en{l} final disparu}
\en{all} \pron{oo} au lieu de \pron{ool}. \en{milk} \pron{mi-k}. \en{people} \pron{piip}.
\textbf{Correction :} le \en{l} final anglais se fait en levant la pointe de la langue vers
les dents du haut, à la fin du mot. Drill : \en{all, call, small, well, tell, feel, milk,
help, people, table, ball, full, cool, school, world, girl}.

\subsection{Erreur 13 --- La liaison à la française là où l'anglais ne lie pas}
En anglais, on ne lie pas un \en{h} (\en{the house} ne devient pas \pron{dhi-zouse}) et on
ne lie pas à travers un mot accentué. En revanche, on lie volontiers consonne \ra voyelle.
\textbf{Correction :} liez seulement quand un mot finit par une consonne et que le suivant
commence par une voyelle : \en{an apple, my old car, turn it on}.

\subsection{Erreur 14 --- Le \en{dark} et le \en{clear} \en{l}, le \en{t} américain}
\en{water} : britannique \pron{ouo-teu} (coup de glotte), américain \pron{oua-deu}.
Choisissez \textbf{une} variété et tenez-vous-y les premiers mois --- mélanger les deux
produit un accent flou. Recommandation : l'américain est plus présent dans les médias,
le britannique est plus enseigné en Afrique de l'Ouest. Les deux sont parfaitement acceptés.

\subsection{Erreur 15 --- Parler trop doucement par peur, ou trop vite par panique}
Trop doucement : on ne vous entend pas, on vous fait répéter, vous perdez confiance.
Trop vite : vous mangez les finales et les accents. \textbf{Correction :} visez un débit
\emph{moyen, régulier}, avec des pauses claires entre les groupes de mots. La pause est un
outil de puissance, pas un aveu de faiblesse : les bons orateurs s'arrêtent.

\section{30 \en{tongue twisters} (virelangues) --- votre salle de sport}
Travaillez-les \textbf{lentement d'abord, puis vite}, trois fois de suite chacun, sans erreur.
Cinq minutes par jour suffisent.

\begin{tableau}[30 virelangues classés par son]
\footnotesize
\begin{longtable}{@{}p{2.6cm}p{7.6cm}p{4.0cm}@{}}
\entete{Son visé} & \entete{Virelangue} & \entete{Traduction} \tabularnewline
\endfirsthead
\entete{Son} & \entete{Virelangue} & \entete{Traduction} \tabularnewline
\endhead
\hline
\ipatheta & \en{The thirty-three thieves thought that they thrilled the throne throughout Thursday.} & Les trente-trois voleurs pensaient avoir ému le trône tout au long de jeudi. \tabularnewline
\ipatheta & \en{Three thin thinkers thought thick thoughts.} & Trois minces penseurs pensaient d'épaisses pensées. \tabularnewline
\ipadh & \en{This is the sixth sick sheikh's sixth sheep.} & C'est la sixième brebis du sixième cheikh malade. \tabularnewline
\ipatheta/\ipadh & \en{They threw three thick things through the thin things.} & Ils ont jeté trois choses épaisses à travers les choses minces. \tabularnewline
/s/ / \ipash & \en{She sells seashells by the seashore.} & Elle vend des coquillages au bord de la mer. \tabularnewline
/s/ / \ipash & \en{Surely Sylvia swims slowly, silently.} & Sûrement Sylvia nage lentement, silencieusement. \tabularnewline
/p/ / b/ & \en{Peter Piper picked a peck of pickled peppers.} & Peter Piper a pris une pincée de poivrons au vinaigre. \tabularnewline
/b/ & \en{A big black bug bit a big black bear.} & Un gros insecte noir a mordu un gros ours noir. \tabularnewline
/f/ / v/ & \en{Four fine fresh figs for you.} & Quatre belles figues fraîches pour toi. \tabularnewline
/w/ & \en{Which witch is which?} & Quelle sorcière est laquelle ? \tabularnewline
/w/ / v/ & \en{Very well, Will we visit the west village Wednesday?} & Très bien, Will, visiterons-nous le village de l'ouest mercredi ? \tabularnewline
/w/ & \en{We surely shall see the sun shine shortly.} & Nous verrons sûrement bientôt briller le soleil. \tabularnewline
/r/ / l/ & \en{Red lorry, yellow lorry.} \quad (×10, vite) & Camion rouge, camion jaune. \tabularnewline
/r/ / l/ & \en{Really leathery, rarely leathery.} & Vraiment coriace, rarement coriace. \tabularnewline
/r/ / l/ & \en{Rolling red wagons, roaring red lions.} & Wagons rouges qui roulent, lions rouges qui rugissent. \tabularnewline
/r/ / w/ & \en{Round the rugged rock the ragged rascal ran.} & Autour du rocher rugueux, le coquin déguenailé courut. \tabularnewline
/l/ & \en{Lovely lemon liniment, lively little lady.} & Joli liniment au citron, petite dame vive. \tabularnewline
/l/ / n/ & \en{Loyal, royal, loyal, royal.} & Loyal, royal, loyal, royal. \tabularnewline
/t/ / d/ & \en{Two tiny tigers take two taxis to town.} & Deux petits tigres prennent deux taxis en ville. \tabularnewline
/d/ & \en{Double Dutch dancers dance daily downtown.} & Les danseurs de Double Dutch dansent tous les jours au centre-ville. \tabularnewline
/k/ / g/ & \en{Good blood, bad blood.} \quad (×10) & Bon sang, mauvais sang. \tabularnewline
/k/ & \en{Can you can a can as a canner can can a can?} & Peux-tu mettre en boîte comme un metteur en boîte peut mettre en boîte ? \tabularnewline
/m/ / n/ & \en{Many mumbling mice are making merry music in the moonlight.} & Beaucoup de souris marmonnantes font de la joyeuse musique au clair de lune. \tabularnewline
/n/ & \en{Nine nice night nurses nursing nicely.} & Neuf gentilles infirmières de nuit soignent gentiment. \tabularnewline
\ipang & \en{The king and the queen are singing in the morning.} & Le roi et la reine chantent le matin. \tabularnewline
\ipae/ \ipaE & \en{A sad man sat on a black cat's mat.} & Un homme triste s'est assis sur le tapis d'un chat noir. \tabularnewline
\ipaee/ \ipaiot & \en{Six sheep sleep in a ship's sheet.} & Six moutons dorment dans le drap d'un navire. \tabularnewline
\ipaou & \en{No one knows how old the old oak grows.} & Personne ne sait quel âge a le vieux chêne. \tabularnewline
/h/ & \en{How much wood would a woodchuck chuck if a woodchuck could chuck wood?} & Combien de bois une marmotte jetterait si une marmotte pouvait jeter du bois ? \tabularnewline
Tout & \en{I thought a thought, but the thought I thought wasn't the thought I thought I thought.} & J'ai pensé une pensée, mais la pensée que j'ai pensée n'était pas celle que je croyais penser. \tabularnewline
\end{longtable}
\end{tableau}

\section{50 phrases-drills : la prononciation en contexte}
Lisez chaque phrase trois fois (lentement, normalement, vite et expressif), puis enregistrez
les 10 premières et comparez à un natif.

\begin{tableau}[50 phrases à maîtriser phonétiquement]
\footnotesize
\begin{longtable}{@{}p{0.7cm}p{7.4cm}p{6.1cm}@{}}
\entete{\#} & \entete{Phrase} & \entete{Point travaillé} \tabularnewline
\endfirsthead
\entete{\#} & \entete{Phrase} & \entete{Point travaillé} \tabularnewline
\endhead
1 & \en{I think this is the third thing they thought about.} & \ipatheta, \ipadh \tabularnewline
2 & \en{He has a house in the heart of the city.} & /h/ aspiré \tabularnewline
3 & \en{We will work with water and wood every week.} & /w/ \tabularnewline
4 & \en{The red car around the world really runs right.} & /r/ anglais \tabularnewline
5 & \en{She should have looked at the book on the shelf.} & finales consonantiques \tabularnewline
6 & \en{The sheep on the ship sleep six sheets deep.} & \ipaee\ contre \ipaiot \tabularnewline
7 & \en{I want a hot cup of tea, not a cold one.} & \ipae, \ipao, \ipalv \tabularnewline
8 & \en{My brother's mother and father live in the north.} & \ipadh, \ipatheta \tabularnewline
9 & \en{No one knows how old the old oak grows.} & diphtongue \ipaou \tabularnewline
10 & \en{The king is singing a long song in the morning.} & \ipang \tabularnewline
11 & \en{It's a beautiful day, isn't it?} & schwa, question tag \tabularnewline
12 & \en{What do you want to do today?} & réductions \en{whaddaya wanna} \tabularnewline
13 & \en{I'm going to have to go now.} & \en{gonna hafta} \tabularnewline
14 & \en{Could you tell him I called?} & \en{couldja tellim} \tabularnewline
15 & \en{I don't know what she's talking about.} & \en{dunno whashe talkinabout} \tabularnewline
16 & \en{Would you like a cup of coffee or a glass of water?} & \en{wouja}, formes faibles \tabularnewline
17 & \en{The weather is better than yesterday.} & \ipadh, comparatif \tabularnewline
18 & \en{He's been working here for about three years.} & \en{present perfect continuous} \tabularnewline
19 & \en{I've never seen anything like that before.} & \en{I've} \pron{aïv}, \en{anything} \tabularnewline
20 & \en{She shouldn't have said that to him.} & \en{shouldn't have} = \pron{chou-dent-ev} \tabularnewline
21 & \en{Where were you when I was waiting for you?} & /w/ + \ipaer\ipalong \tabularnewline
22 & \en{My father's favourite film is a French film.} & \ipaer\ipalong, /f/ \tabularnewline
23 & \en{Actually, I'd rather have a vegetable salad.} & réductions de syllabes \tabularnewline
24 & \en{The library is behind the bakery, near the river.} & \en{library} \pron{LAÏ-bre-ri} \tabularnewline
25 & \en{February is the shortest month of the year.} & \en{February} \pron{FÉ-brou-é-ri} \tabularnewline
26 & \en{I usually go to the university by bus.} & /\ipazh/ dans \en{usually} \tabularnewline
27 & \en{The children are playing in the garden.} & \en{children} \pron{TCHIL-dren} \tabularnewline
28 & \en{Question after question, he never answered.} & \en{question} \pron{KWÈS-tchoun} \tabularnewline
29 & \en{She's a photographer and he's an engineer.} & accents mobiles \tabularnewline
30 & \en{It's not very interesting, to be honest.} & \en{honest} sans /h/ \tabularnewline
31 & \en{Half of the class was absent, which is unusual.} & \en{half} sans /l/ \tabularnewline
32 & \en{Could I have a knife and a fork, please?} & \en{knife} sans /k/ \tabularnewline
33 & \en{He wrote a wrong address on the envelope.} & \en{wrote, wrong} sans /w/ \tabularnewline
34 & \en{Climbing the mountain was hard but worth it.} & \en{climb} sans /b/, \en{worth} \tabularnewline
35 & \en{I'm looking forward to hearing from you.} & bloc social essentiel \tabularnewline
36 & \en{Would you mind opening the window?} & politesse + \en{would you} \tabularnewline
37 & \en{I'm not sure whether he'll come or not.} & \en{whether} \pron{WÈ-dheu} \tabularnewline
38 & \en{The problem is that we don't have enough time.} & \en{enough} \pron{i-NAF} \tabularnewline
39 & \en{She's used to working late, so it's fine.} & \en{used to} \pron{iou-s-tou} \tabularnewline
40 & \en{What would you have done in my place?} & conditionnel 3, \en{would you} \tabularnewline
41 & \en{I've been living here since I was a child.} & \en{since} + voyelles longues \tabularnewline
42 & \en{There's a lot of traffic at this hour.} & \en{hour} sans /h/, \en{a lot of} \tabularnewline
43 & \en{Although it's cold, I'm going out for a walk.} & \en{although} \pron{o-l-DZOOU} \tabularnewline
44 & \en{He asked me whether I had heard the news.} & \en{asked} \pron{aaskt}, discours rapporté \tabularnewline
45 & \en{You should have told me earlier, honestly.} & \en{should have} + \en{honestly} \tabularnewline
46 & \en{Everybody knows that nobody knows everything.} & \en{everybody} \pron{È-vri-ba-di} \tabularnewline
47 & \en{I'd like a glass of water without ice, please.} & \en{water}, \en{without} \tabularnewline
48 & \en{The government has decided to raise taxes.} & \en{government} \pron{GOU-veun-m'nt} \tabularnewline
49 & \en{It's probably going to rain this afternoon.} & \en{probably} \pron{PROB-li} \tabularnewline
50 & \en{Thank you very much for your help, I really appreciate it.} & remerciement complet \tabularnewline
\end{longtable}
\end{tableau}

\section{Programme de correction en 30 jours}
\begin{tableau}[30 jours pour réparer votre accent]
\begin{tabularx}{\linewidth}{@{}p{2.0cm}X@{}}
\entete{Jours} & \entete{Travail quotidien (10 à 15 minutes)} \tabularnewline
\hline
\textbf{1 à 5} & \ipatheta\ et \ipadh\ : 40 mots par jour devant un miroir + 3 virelangues. \tabularnewline
\textbf{6 à 9} & /h/ et /w/ : 40 mots + paires minimales \en{west/vest} + 3 virelangues. \tabularnewline
\textbf{10 à 13} & /r/ et /l/ : 40 mots + \en{red lorry yellow lorry} ×10 + lecture à voix haute de 5 phrases-drills. \tabularnewline
\textbf{14 à 17} & Voyelles : paires \en{sheep/ship}, \en{full/fool}, \en{cart/cut}, \en{cat/cut}, \en{no/now} + drill des 50 phrases (1 à 15). \tabularnewline
\textbf{18 à 21} & Finales consonantiques et \ipang\ : aucun « euh » ajouté + drill des 50 phrases (16 à 30). \tabularnewline
\textbf{22 à 25} & Rythme et accents : shadowing quotidien d'un dialogue de ce livre + marquage des accents au crayon. \tabularnewline
\textbf{26 à 28} & \en{Connected speech} : réductions (\en{gonna, wanna, gotta}) et formes faibles + drill des 50 phrases (31 à 50). \tabularnewline
\textbf{29 à 30} & \textbf{Bilan :} enregistrez 5 minutes de parole libre, comparez avec votre enregistrement du jour 1, listez vos 3 sons restants à travailler en continu. \tabularnewline
\end{tabularx}
\end{tableau}

\begin{memo}[Ce qu'il faut retenir de la Partie II]
\begin{enumerate}
  \item L'anglais a \textbf{9 sons inexistants en français} : \ipatheta, \ipadh, /h/, /w/,
  /r/, \ipang, \ipae, \ipalv, et le schwa. Travaillez-les un par un, dans cet ordre.
  \item \textbf{L'accent tonique} est plus important que les sons : \en{baNAna}, pas
  \en{banana}.
  \item Le \textbf{schwa} est la clé du rythme : écrasez les syllabes non accentuées.
  \item Le \textbf{\en{connected speech}} explique 80\,\% de vos difficultés d'écoute :
  les natifs ne prononcent pas les mots séparément.
  \item Choisissez \textbf{une} variété (américaine ou britannique) pour les 6 premiers mois.
  \item Enregistrez-vous \textbf{chaque semaine} : c'est le seul miroir honnête.
\end{enumerate}
\end{memo}
